\subsection{Decomposition into slices. Inverse image of a measure under a local homeomorphism}

Let X be a locally compact space, \(\pi\) a mapping of X into a locally compact space T, \(\mu\) a positive measure on T, \(\Lambda: t \mapsto \lambda_t\) a scalarly essentially \(\mu\) -integrable and vaguely \(\mu\) -measurable mapping of T into \(\mathcal{M}_{+}(X)\) . Let \(\nu = \int \lambda_t d\mu(t)\) . If \(\lambda_t\) is carried by \(\overline{\pi}(t)\) for every \(t \in T\) , the equality \(\nu = \int \lambda_t d\mu(t)\) is said to be a decomposition into slices (or a disintegration) of \(\nu\) relative to \(\pi\) . This concept will be studied in detail in Ch. VI.

PROPOSITION 10. --- With the above notations, suppose that \(\pi\) is \(\nu\) -measurable. Let \(g\) be the function \(t \mapsto \lambda_t^*(1)\) on T. For \(\pi\) to be \(\nu\) -proper, it is necessary and sufficient that \(g\) be locally \(\mu\) -integrable, in which case

\[
\pi (\nu) = g \cdot \mu .\tag{14}
\]

We begin by arguing on the assumption that g is finite locally \(\mu\) -almost everywhere; we will rid ourselves of this auxiliary hypothesis at the end of the proof. Since \(\pi\) is by hypothesis \(\nu\) -measurable, to say that \(\pi\) is \(\nu\) -proper is equivalent to saying that \(\nu^{\bullet}(f\circ\pi)<+\infty\) for every function \(f\in\mathcal{K}_{+}(\mathrm{T})\) ; g being finite locally almost everywhere, we are under the conditions for applying assertion c) of Prop. 5 of §3, No.~2. Therefore

\[
\int^ {\bullet} (f \circ \pi) d \nu = \int^ {\bullet} d \mu (t) \int^ {\bullet} (f \circ \pi) d \lambda_ {t} = \int^ {\bullet} f (t) g (t) d \mu (t),
\]

from the fact that \(\lambda_{t}\) is concentrated on \(\overline{\pi}^{1}(t)\) . We know that g is \(\mu\) -measurable, since \(\Lambda\) is \(\mu\) -adequate ( \(\S3\) , No.~1, Def. 1). To say that the first member is finite for every \(f \in \mathcal{K}_{+}(\mathrm{T})\) is therefore equivalent to saying that g is locally \(\mu\) -integrable ( \(\S5\) , Prop. 1), and in this case (14) follows at once from the above relations.

It therefore only remains to eliminate the auxiliary hypothesis. If g is locally \(\mu\) -integrable, then g is finite locally \(\mu\) -almost everywhere, and the hypothesis is indeed satisfied. Let us assume that \(\pi\) is \(\nu\) -proper, and let us show that g is finite locally almost everywhere. Let \(\mathfrak{K}\) be the \(\mu\) -dense set of compact sets K such that \(\Lambda|K\) is vaguely continuous; since g is measurable, we are reduced to showing that every compact set \(K \in \mathfrak{K}\) such that \(g|K = +\infty\) is \(\mu\) -negligible. Now, let \(\mathcal{H}\) be the set of functions \(h \in \mathcal{K}_{+}(X)\) such that \(h \leqslant 1\) ; set \(g_{h}(t) = \lambda_{t}(h)\) , denote by \(\Lambda_{h}\) the \(\mu\) -adequate mapping \(t \mapsto h \cdot \lambda_{t}\) , by \(\nu_{h}\) the integral of \(\Lambda_{h}\) , and by f an element of \(\mathcal{K}_{+}(T)\) such that \(f \geqslant \varphi_{K}\) . Applying formula (14) to \(\Lambda_{h}\) , which does satisfy the auxiliary hypothesis, we obtain:

\[
\int (f \circ \pi) d \nu \geqslant \int (f \circ \pi) d \nu_ {h} = \int f g _ {h} d \mu .
\]

But the functions \(fg_h|K\) form an increasing directed set of continuous functions on \(K\) , whose upper envelope has the value \(+\infty\) ; by Dini's theorem (GT, X, §4, No.~1, Th. 1), one can choose \(h\) so that \(fg_h|K\) is greater than or equal to an arbitrary positive number \(n\) , and it follows that \(\int(f \circ \pi) d\nu \geqslant n \mu(K)\) . Since the first member is finite because \(\pi\) is \(\nu\) -proper, it then follows that \(\mu(K) = 0\) .

COROLLARY 1. --- Suppose that \(\pi\) is \(\nu\) -measurable.

\begin{enumerate}
\def\labelenumi{\alph{enumi})}
\item
  If \(\mathrm{N}\subset \mathrm{T}\) is locally \(\mu\) -negligible, then \(\overline{\pi}^{1}(\mathrm{N})\) is locally \(\nu\) -negligible.
\item
  If \(f\) is a \(\mu\) -measurable mapping of \(T\) into a topological space \(G\) , then \(f \circ \pi\) is \(\nu\) -measurable.
\end{enumerate}

We take up again the notations \(\Lambda_h\) , \(\nu_h\) , \(g_h\) of the end of the preceding proof: \(\nu_h\) being a bounded measure for every \(h \in \mathcal{H}\) , \(\pi\) is \(\nu_h\) -proper, \(g_h\) is locally \(\mu\) -integrable, and \(\pi(\nu_h) = g_h \cdot \mu\) , a measure with base \(\mu\) . It follows that N is locally negligible (resp. that \(f\) is measurable) for the measure \(\pi(\nu_h)\) ( \(\S 5\) , No.~3, Cor. 1 of Prop. 3 and Prop. 4). Consequently \(\overline{\pi}^1(N)\) is locally negligible (resp. \(f \circ \pi\) is measurable) for the measure \(\nu_h\) (Cor. 2 of

Prop. 2, resp. Prop. 3). Finally, one notes that the measures \(\nu_{h}\) form an increasing directed family of positive measures whose supremum is \(\nu\) , and one applies Cor. 1 (resp. Cor. 2) of Prop. 11 of §1, No.~4.

COROLLARY 2. --- Suppose that \(\pi\) is \(\nu\) -proper; let \(\mathbf{f}\) be a mapping defined on \(\mathbf{T}\) , with values in a Banach space or in \(\overline{\mathbf{R}}\) . For \(f \circ \pi\) to be essentially \(\nu\) -integrable, it is necessary and sufficient that \(gf\) be essentially \(\mu\) -integrable.

Taking into account Prop. 10, this follows immediately from Th. 1 of §5, No.~3 and Th. 1 of No.~2.

Example. --- Let X and T be two locally compact spaces, and let \(\pi\) be a local homeomorphism of X into T. In other words (GT, I, §11, Exer. 25), we assume that every point \(x \in X\) admits a neighborhood V such that \(\pi|V\) is a homeomorphism of V onto a neighborhood of \(\pi(x)\) ; if necessary replacing V by a relatively compact open neighborhood W of x such that \(\overline{W} \subset V\) , one deduces that the set \(\mathcal{U}\) of relatively compact open subsets U of X, such that \(\pi|\overline{U}\) is a homeomorphism of \(\overline{U}\) onto its image, is an open covering of X. Now let \(\mu\) be a positive measure on T; if U is an element of \(\mathcal{U}\), then \(\pi(U)\) is an open set in the compact space \(\pi(\overline{U})\) , therefore is a locally compact subspace of T, and one knows how to define the measure \(\mu|\pi(U)\) induced by \(\mu\) on \(\pi(U)\) (Ch. IV, §5, No.~7). Let \(\nu_{U}\) be the image of \(\mu|\pi(U)\) under the homeomorphism inverse to \(\pi|U\) ; we are going to show that there exists one and only one measure \(\nu\) on X that induces the measure \(\nu_{U}\) on every open set \(U \in \mathcal{U}\) . This measure is called the inverse image of \(\mu\) under the local homeomorphism \(\pi\) , and is denoted \(\overline{\pi}^{1}(\mu)\) .

The uniqueness of \(\nu\) follows at once from the principle of localization (Ch. III, §2, No.~1, Cor. of Prop. 1). To establish existence, we note that if \(t \in \mathbf{T}\) , then every point \(x \in \overline{\pi}^{-1}(t)\) admits a neighborhood that intersects \(\overline{\pi}^{-1}(t)\) only at the point \(x\) , so that \(\overline{\pi}^{-1}(t)\) is a discrete subspace of \(\mathbf{X}\) , and that the family \((\varepsilon_x)_{x \in \overline{\pi}^{-1}(t)}\) is summable; we denote its sum by \(\lambda_t\) . We next show that the mapping \(t \mapsto \lambda_t\) is scalarly essentially \(\mu\) -integrable, and that its integral \(\nu = \int \lambda_t d\mu(t)\) is the sought-for inverse image. This will result at once from the following lemma:

Lemma. --- a) Let \(f\) be an element of \(\mathcal{K}_{+}(X)\) ; the function \(t \mapsto \lambda_{t}(f)\) is positive, upper semi-continuous, with compact support, and its restriction to \(\pi(X)\) is continuous.

\begin{enumerate}
\def\labelenumi{\alph{enumi})}
\setcounter{enumi}{1}
\tightlist
\item
  Let U be an element of \(\mathcal{U}\), \(\nu\) the integral of the scalarly essentially \(\mu\) -integrable function \(t \mapsto \lambda_{t}\) ; the image of the measure \(\nu|U\) under \(\pi|U\) is equal to \(\mu|\pi(U)\) .
\end{enumerate}

To establish a), one may reduce by means of a partition of unity (Ch. III, §1, No.~2, Lemma 1) to the case that the support S of f is contained in an open set \(U \in \mathcal{U}\) . Let g be the mapping \(t \mapsto \lambda_{t}(f)\) ; \(\pi|U\) being a homeomorphism, \(g|\pi(U)\) belongs to \(\mathcal{K}_{+}(\pi(U))\) , consequently \((\pi(U)\) being an open set in \(\pi(X)\) ) the restriction of g to \(\pi(X)\) is continuous. Since g is positive and the restriction of g to the compact set \(\pi(S)\) is continuous, one sees that g is upper semi-continuous on T. It follows that g is \(\mu\) -integrable.

To establish b), denote by g an element of \(\mathcal{K}\big(\pi(\mathrm{U})\big)\) , by \(g^{\circ}\) its extension by 0 to T, by f the function \(g \circ (\pi|U)\) , and by \(f^{\circ}\) the extension by 0 of f to X. The assertion b) is equivalent to the equality \(\int g^{\circ} d\mu = \int f^{\circ} d\nu\) . But \(f \in \mathcal{K}(\mathrm{U})\) , therefore \(f^{\circ} \in \mathcal{K}(\mathrm{X})\) , and the second integral is therefore equal to \(\int \lambda_{t}(f^{\circ}) d\mu(t)\) . Finally \(\lambda_{t}(f^{\circ}) = g^{\circ}(t)\) , which completes the proof.

We now observe that \(\pi(\mathrm{X})\) is open in T, hence is \(\mu\) -measurable; the mapping \(\Lambda: t \mapsto \lambda_t\) is vaguely \(\mu\) -measurable, because its restriction to each of the sets \(\pi(\mathrm{X})\) and \(\mathbb{C}\pi(\mathrm{X})\) is vaguely continuous. Under these conditions, the formula \(\overline{\pi}^{1}(\mu) = \int \lambda_t d\mu(t)\) defines a decomposition into slices of \(\overline{\pi}^{1}(\mu)\) relative to \(\pi\) , and Prop. 10 yields the following result:

PROPOSITION 11. --- Let \(\pi\) be a local homeomorphism of a locally compact space \(X\) into a locally compact space \(T\) , and let \(\mu\) be a positive measure on \(T\) . Let \(n\) be the numerical function that associates to every \(t \in T\) the number of elements of \(\overline{\pi}^{1}(t)\) if this number is finite, and \(+\infty\) in the contrary case. For \(\pi\) to be \(\overline{\pi}^{1}(\mu)\) -proper, it is necessary and sufficient that \(n\) be locally \(\mu\) -integrable, in which case

\[
\pi \big (\bar {\pi} ^ {1} (\mu) \big) = n \cdot \mu .\tag{15}
\]

